% Generated draft: requires mathematical and source review.
\documentclass[11pt]{article}
\usepackage[margin=25mm]{geometry}
\usepackage{fontspec}
\setmainfont{Latin Modern Roman}
\usepackage{amsmath,amssymb,mathtools,unicode-math}
\setmathfont{Latin Modern Math}
\usepackage{enumitem,needspace,xurl,hyperref}
\setlength{\parindent}{0pt}
\setlength{\parskip}{5pt}
\setlength{\emergencystretch}{3em}
\newcommand{\N}{\mathbb{N}}
\newcommand{\Z}{\mathbb{Z}}
\newcommand{\Q}{\mathbb{Q}}
\newcommand{\R}{\mathbb{R}}
\newcommand{\C}{\mathbb{C}}
\newcommand{\F}{\mathbb{F}}
\newcommand{\A}{\mathbb{A}}
\newcommand{\PP}{\mathbb{P}}
\newcommand{\E}{\mathbb{E}}
\newcommand{\eps}{\varepsilon}
\newcommand{\dd}{\,\mathrm{d}}
\begin{document}
\section*{UnsolvedMath ID: 30000054}
\textbf{Source problem code:} catalogue entry
\subsection{Definitions and mathematical statement}
The source asks to compare recount risk under Electoral College and national popular-vote aggregation, but does not specify a probability model or recount rule. A precise comparison must use one declared joint vote model and two declared event definitions.

For example, let \(M_i\) be the net two-candidate vote count and \(V_i\) turnout in state \(i\), with joint law \(\Pi\). The national net count is then \(M_{\mathrm{nat}}=\sum_iM_i\); a fractional national margin instead uses \(\sum_iM_i/\sum_iV_i\). Electoral weights, allocation rules, and tie rules determine the Electoral College outcome. The source does not supply these data, nor independence, dependence, turnout, or candidate-conditioning assumptions.

After choosing recount thresholds and rules, let \(E_{\mathrm{EC}}\) be the event called a recount-triggering Electoral College victory and \(E_{\mathrm{PV}}\) the event called a recount-triggering national popular-vote margin. The requested comparison is
\[
 \Pi(E_{\mathrm{EC}})>\Pi(E_{\mathrm{PV}}).
\]
A common popular-vote rule might use \(|M_{\mathrm{nat}}|\leq r_{\mathrm{nat}}\), but the source does not choose count versus percentage thresholds. For the Electoral College event it must additionally specify whether any close state, a close state pivotal to the winner, or a collection of jointly pivotal states triggers the event. The same unconditional or victory-conditioned convention must be used in both probabilities. None of these possible completions is silently adopted as the source's intended one.

If an asymptotic comparison is intended, the number of voters/states, joint-law sequence, and scaling of recount thresholds must also be fixed. The record supplies no such limiting regime and no quantifier identifying which model class should satisfy the inequality. It therefore does not define a universal probabilistic theorem.

The supplied triage dated 21 August 2026 reports that model-based studies support more frequent narrow, disputable outcomes under Electoral College aggregation, while explicitly calling the displayed comparison underspecified. That empirical/model-dependent assessment is not a proof of the exact recount-event inequality for every joint vote law or threshold choice.
\subsection{Short English statement}
The comparison is \(\Pi(E_{\mathrm{EC}})>\Pi(E_{\mathrm{PV}})\), but the joint law, thresholds, pivotality rule, and conditioning are missing. The supplied literature assessment is model-dependent and does not fix a universal theorem or asymptotic regime.
\subsection{Sources}
\begin{itemize}
\item Local UnsolvedMath registry, permanent ID \(30000054\), exact source code \texttt{catalogue entry}, Oberwolfach Reports Open Problems. Statement status \texttt{exact}; verification reports coherent wording needing no textual repair.
\item ``Analysis and Design of Electoral Systems (2004),'' Oberwolfach Reports, pp. 723--772, DOI \texttt{10.4171/owr/2004/14}: \url{https://doi.org/10.4171/owr/2004/14}. Proposed year 2004.
\item Supplied ``Geruso--Spears, Electoral College risk'': \url{https://www.nber.org/papers/w27993}; supplied ``OWR 14/2004'': \url{https://doi.org/10.4171/OWR/2004/14}. Fuller bibliographic entries and the exact compared recount events are not supplied.
\item Literature checked 21 August 2026: supplied status \texttt{open}, with explicit joint-distribution and threshold-scaling limitations. Registry created 25 August 2026 and updated 16 September 2026. This dated curation triage requires independent verification before citation.
\item Registry curation attribution: Ulam AI / UnsolvedMath Contributors, CC BY 4.0. Underlying report and literature retain their supplied provenance; this statement-only draft requires human review.
\end{itemize}
\end{document}
